\SourceLeaf{16}
\par\medskip\Needspace{4\baselineskip}
\centerline{THEOREMA SEXTUM.}\medskip
\par
{\itshape Attractio corporis in punctum $M$, axi coordinatarum $x$ parallela et opposita, exhibetur per integrale
\[
-\int\frac{\mathrm{d}s.\Phi r.\cos MQ.\cos MX}{rr}
\]
per totam superficiem extensum, sive $M$ iaceat extra corpus, sive intra, sed adiecta parte $-\pi\Phi0.\cos QX$, quoties $M$ iacet in ipsa superficie, ubi pro $QX$ accipiendus est valor definitus, quem habet in $M$.}\par
Manifesto vires secundum directiones axibus coordinatarum $y,z$ parallelas atque oppositas perinde exprimentur per integralia
\[
-\int\frac{\mathrm{d}s.\Phi r.\cos MQ.\cos MY}{rr},\quad-\int\frac{\mathrm{d}s.\Phi r.\cos MQ.\cos MZ}{rr}
\]
quibus adiicere oportet $-\pi\Phi0.\cos QY$, $-\pi\Phi0.\cos QZ$ (sumtis pro angulis valoribus definitis in $M$), quoties $M$ iacet in corporis superficie.

Ceterum facile perspicitur, tres vires
\[
-\pi\Phi0.\cos QX,\quad-\pi\Phi0.\cos QY,\quad-\pi\Phi0.\cos QZ
\]
aequivalere unicae $=-\pi\Phi0$ ipsi superficiei normali introrsumque directae.

Manifesto evolutione integralis $-\Phi0.\int\mathrm{d}\Sigma.\cos MX$ supersedere potuissemus, si functio $f$ ita comparata est. ut liceat statuere $\Phi0=0$; sed maluimus disquisitionem omni generalitate persequi. Quoties autem attractio cubo altiorive potestati distantiae inverse proportionalis supponitur, patet, illud non licere, sed necessario fieri $\Phi0=-\infty$, unde sequitur, in tali suppositione punctum in corporis superficie positum vi infinita versus solidum premi.

\Article{9}
\par
Per methodos hactenus explicatas integralia, quae per totum corporis volumen extendi debuissent (integralia tripla). ad talia reduximus, quae tantummodo per corporis superficiem sunt extendenda, et quidem duplici modo. Indoles superficiei exprimitur per aequationem inter coordinatas $x,y,z$, i. e. per aequationem $W=0$, denotante $W$ functionem variabilium $x,y,z$, quam ab omni irrationalitate liberam supponere licet. Prodeat e differentiatione functionis $W$
\[
\mathrm{d}W=T\mathrm{d}x+U\mathrm{d}y+V\mathrm{d}z
\]
